\subsection{Adaptive critical experiments}
\label{sec:v16-adaptive-critical}
The exact overlap calculation also has a stopped, adaptive form. The
relevant overlap mass is generally not a product of unconditional
one-step overlaps, because later windows depend on earlier records.
We give its exact expression and a physical approximation bound.
Throughout this subsection the comparison is the same simple binary
pair as in Section~\ref{sec:v7-critical}: a fixed table's finite law
versus its boundary model. The policy is independent of the binary
hypothesis, but may use the specified table and the past records.

For a design $a=(j,d)$ write $P_a^0=P_{j,d}^0$ and
$Q_a^0=P_{\partial,j,d}^0$. Define
\begin{equation}\label{eq:v16-overlap-hazard}
 h(a)=p_{j,d}^0\delta_j,
 \qquad \delta_j=\frac2\pi\arcsin(e^{-j\gamma}),\qquad
 p_{j,d}^0=\frac{d^2}{2A\sinh(j\gamma)}.
\end{equation}
Their densities agree on the common support, including the failure
atom. Let $C_a$ be their common restriction and put
$R_a=C_a/(1-h(a))$. The chosen collar ensures $h(a)<1$.
For a policy $\pi$ with cap $N$, let $\mathsf R^\pi$ be the law
obtained by following the same policy with record kernel $R_a$.
This auxiliary law is supported on histories that have not revealed
the binary hypothesis. Its stopping rule is still the specified rule
of $\pi$.

\begin{theorem}[Exact stopped overlap and physical transfer]
\label{thm:v16-adaptive-critical}
The tangent stopped transcript laws satisfy
\begin{equation}\label{eq:v16-overlap-functional}
 \|\mathsf P^{0,\pi}-\mathsf Q^{0,\pi}\|_{\rm TV}
     =1-m_\pi,\qquad
 m_\pi=\mathbb E_{\mathsf R^\pi}
                   \prod_{i=1}^{T}(1-h(a_i)).
\end{equation}
As a simple binary experiment, this pair is exactly equivalent by
hypothesis-independent randomizations to an erasure experiment with
erasure probability $m_\pi$. Its minimum equal-prior testing error is
$m_\pi/2$.

Let $\mathsf P^\pi,\mathsf Q^\pi$ instead be the actual positive-offset
finite and boundary transcript laws. All reachable designs are assumed
to lie in the common collar. Put
\begin{equation}\label{eq:v16-adaptive-tangent-budget}
 \varepsilon_\pi=
   \mathbb E_{\mathsf P^\pi}\sum_{i=1}^{T}p_{a_i}^0\omega(d_i)
  +\mathbb E_{\mathsf Q^\pi}\sum_{i=1}^{T}p_{a_i}^0\omega(d_i),
\end{equation}
where $\omega(d)=\sqrt d$, or $\omega(d)=d$ when both contact graphs
are even. Uniformly on the compact geometric families under consideration,
\begin{equation}\label{eq:v16-adaptive-critical-physical}
 \left|\|\mathsf P^\pi-\mathsf Q^\pi\|_{\rm TV}
                                 -(1-m_\pi)\right|
                           \le C\varepsilon_\pi.
\end{equation}
The randomizations in the first assertion concern the specified known
simple pair and policy, not an unknown family of billiard tables.
\end{theorem}
\begin{proof}
At each design, the measures $P_a^0,Q_a^0$ have the same restriction
$C_a$ to their overlap; their other restrictions are supported on disjoint
sets. A transcript compatible with both hypotheses can therefore contain
only common-support records at every step. Conversely, along any such
history every observation density agrees under the two hypotheses.
The design kernels and stopping probabilities agree as well, since
they are functions of the same observed past. Thus the two full transcript
densities agree on their overlap. A first exclusive record makes the
entire transcript exclusive, regardless of all later decisions.

The common transcript restriction is obtained by using the
subprobability kernels $C_a$ at active steps and the original policy
kernels for choosing designs and stopping. Since
$C_a=(1-h(a))R_a$, its density relative to $\mathsf R^\pi$ is
\[
                         Z=\prod_{i=1}^{T}(1-h(a_i)).
\]
This identity follows by induction over the at most $N$ active steps,
keeping the stopping branch at each step; it also covers $T=0$ by the
empty-product convention. Its total mass is $m_\pi$. Equal densities
on the overlap give total variation $1-m_\pi$.

Map the overlap to an erasure symbol and each exclusive part to its
hypothesis label. For the reverse randomization, given an erasure draw
from the normalized common transcript restriction; given a label draw
from that hypothesis's normalized exclusive restriction. The common
restriction is $Z\,\mathsf R^\pi$, not in general $\mathsf R^\pi$
itself. Assign any fixed distribution to a zero-mass branch. These
kernels recover both laws and do not depend on the true binary label.
With equal priors the posterior remains equal on an erasure and the
hypothesis is known otherwise, proving the testing assertion.

By Lemma~\ref{lem:v7-tangent}, the actual raw kernel and its own
tangent kernel have total variation at most $Cp_a^0\omega(d)$, for
either hypothesis. Apply Lemma~\ref{lem:v16-kernel-comparison}
separately to these two approximations, under the same policy.
The sum of the resulting bounds is $C\varepsilon_\pi$. The reverse
triangle inequality for total variation then proves
\eqref{eq:v16-adaptive-critical-physical}.
\end{proof}

\begin{corollary}[An adaptive critical limit]
\label{cor:v16-adaptive-critical-limit}
Consider a sequence of the preceding experiments and capped policies.
Under the corresponding auxiliary laws $\mathsf R^{\pi_n}$, suppose
\begin{equation}\label{eq:v16-adaptive-hazard-limit}
 \sum_{i=1}^{T_n}h(a_i)\longrightarrow b\in[0,\infty],\qquad
 \max_{i\le T_n}h(a_i)\longrightarrow0
 \quad\hbox{in probability},
\end{equation}
where the maximum of an empty list is zero. If
$\varepsilon_{\pi_n}\to0$, the actual stopped transcript distances
converge to $1-e^{-b}$ and their minimum equal-prior testing errors
converge to $e^{-b}/2$. For $b=\infty$, convergence of the sum to
infinity and $\varepsilon_{\pi_n}\to0$ suffice without the maximum
condition.
\end{corollary}
\begin{proof}
For $0\le x<1$, the series of $-\log(1-x)$ gives
\[
 0\le-\log(1-x)-x\le\frac{x^2}{2(1-x)}.
\]
For finite $b$, the sum of squared hazards is at most their maximum
times their sum, hence tends to zero in probability. Applying the
displayed inequality on the event that the maximum is at most $1/2$
shows that the product in \eqref{eq:v16-overlap-functional} converges
in probability to $e^{-b}$. Products lie in $[0,1]$, so their
expectations converge as well: for each $\epsilon>0$, the expectation
of their absolute error is at most $\epsilon$ plus the probability
that the error exceeds $\epsilon$. If $b=\infty$, bound the product
by the exponential of minus the hazard sum and use the same boundedness
argument. Finally apply \eqref{eq:v16-adaptive-critical-physical} and
the binary testing identity.
\end{proof}

For predetermined windows and a deterministic stopping index,
\eqref{eq:v16-overlap-functional} reduces to
\eqref{eq:v7-raw-product}. For adaptive windows it retains the observed
common-history distribution instead of incorrectly averaging each hazard
in isolation. The conditions concern the specified binary comparison;
they do not assert optimal design or a minimax lower bound for geometric
reconstruction. All transcript laws in this subsection retain the failure
atom and the actual number of preparations.
